% Part: incompleteness
% Chapter: incompleteness-provability
% Section: first-incompleteness-thm

\documentclass[../../../include/open-logic-section]{subfiles}

\begin{document}

\olfileid{inc}{inp}{1in}

\olsection{প্রথম অসম্পূর্ণতা উপপাদ্য}

এবার প্রথম অসম্পূর্ণতা উপপাদ্যের গ্যোডেলীয় মূল প্রমাণটি বর্ণনা করতে
পারি। ধরি $\Th{T}$ হলো পাটিগণিতের ভাষাকে প্রসারিত করে এমন কোনো ভাষায়
গণনসাধ্যভাবে স্বতঃসিদ্ধায়িত তত্ত্ব, এবং $\Th{T}$-তে $\Th{Q}$-এর
স্বতঃসিদ্ধগুলি আছে। এর বিশেষ অর্থ হলো, $\Th{T}$ গণনসাধ্য অপেক্ষক ও
সম্বন্ধগুলিকে উপস্থাপন করে।

আমরা যুক্তি দিয়েছি যে, সূত্র ও প্রমাণগুলিকে সংখ্যা দিয়ে যুক্তিসঙ্গতভাবে
সংকেতায়িত করলে $\Prf[\Th{T}](x,y)$ সম্বন্ধটি গণনসাধ্য; এখানে
$\Prf[\Th{T}](x,y)$ তখনই সত্য, যখন $x$ হলো $\Th{T}$-তে গ্যোডেল
সংখ্যা~$y$-বিশিষ্ট !!{formula}-টির !!a{derivation}-এর গ্যোডেল সংখ্যা।
বস্তুত গ্যোডেলের বিবেচিত বিশেষ তত্ত্বটির ক্ষেত্রে তাঁর গবেষণাপত্রের ৪৫টি
অপেক্ষক ও সম্বন্ধের তালিকা ব্যবহার করে তিনি দেখাতে পেরেছিলেন যে এই
সম্বন্ধটি আদিম পুনরাবৃত্ত। তাঁর নির্দিষ্ট $\Th{T}$-এর জন্য তালিকার ৪৫তম
সম্বন্ধ $x B y$ ঠিক $\Prf[\Th{T}](x,y)$। মনে রেখো, গ্যোডেল তাঁর
গবেষণাপত্রে যেখানে ``পুনরাবৃত্ত'' শব্দটি ব্যবহার করেন, সেখানে এখন আমরা
``আদিম পুনরাবৃত্ত'' বলি।

$\Prf[\Th{T}](x,y)$ গণনসাধ্য বলে সেটি $\Th{T}$-তে উপস্থাপনযোগ্য। তাকে
উপস্থাপনকারী সূত্রটি $\OPrf[\Th{T}](x,y)$ দিয়ে বোঝাব।
$\OProv[\Th{T}](y)$ হলো $\lexists[x][\OPrf[\Th{T}](x,y)]$ সূত্রটি।
এটি গ্যোডেলের তালিকার ৪৬তম সম্বন্ধ $\fn{Bew}(y)$-কে বর্ণনা করে। গ্যোডেল
যেমন বলেন, এটিই একমাত্র সম্বন্ধ যাকে ``পুনরাবৃত্ত বলা যায় না।'' সম্ভবত
তিনি বোঝাতে চেয়েছিলেন: সংজ্ঞা থেকে এটি গণনসাধ্য কি না স্পষ্ট নয়; পরের
বিকাশগুলি বস্তুত দেখায় যে এটি গণনসাধ্য নয়।

ধরি $\Th{T}$ হলো $\Th{Q}$-কে অন্তর্ভুক্ত করে এমন একটি
!!{axiomatizable} তত্ত্ব। তখন $\Prf[\Th{T}](x, y)$ নির্ণেয়, অতএব
!!a{formula}~$\OPrf[\Th{T}](x, y)$ দ্বারা $\Th{Q}$-তে উপস্থাপনযোগ্য।
$\OProv[\Th{T}](y)$ আগের বর্ণিত সূত্র। স্থির-বিন্দু লেমা অনুসারে এমন একটি
সূত্র $!G_\Th{T}$ আছে, যাতে $\Th{Q}$ (এবং তাই $\Th{T}$) নিচেরটি
!!{derive}s করে:
\begin{equation}
\ollabel{eqn:qpf}
!G_\Th{T} \liff \lnot \OProv[\Th{T}](\gn{!G_\Th{T}}).
\end{equation}
খেয়াল করো, $!G_\Th{T}$-এর সারকথা হলো, ``$!G_\Th{T}$ তত্ত্ব
$\Th{T}$-তে !!{derivable} নয়।''

\begin{lem}\ollabel{lem:cons-G-unprov}
$\Th{T}$ যদি $\Th{Q}$-কে প্রসারিত করে এমন একটি সঙ্গতিপূর্ণ,
!!{axiomatizable} তত্ত্ব হয়, তবে $\Th{T} \Proves/ !G_\Th{T}$।
\end{lem}

\begin{proof}
ধরি $\Th{T}$ সূত্র $!G_\Th{T}$-কে !!{derive}s করে। তবে সত্যিই একটি
!!a{derivation} আছে; ফলে কোনো সংখ্যা $m$-এর জন্য
$\Prf[\Th{T}](m, \Gn{!G_\Th{T}})$ সত্য। কিন্তু তখন $\Th{Q}$ বাক্য
$\OPrf[\Th{T}](\num m, \gn{!G_\Th{T}})$-কে !!{derive}s করে। অতএব
$\Th{Q}$ $\lexists[x][\OPrf[\Th{T}](x,\gn{!G_\Th{T}})]$-কে !!{derive}s
করে, যা সংজ্ঞা অনুসারে $\OProv[\Th{T}](\gn{!G_\Th{T}})$।
\olref{eqn:qpf} অনুসারে $\Th{Q}$ সূত্র $\lnot !G_\Th{T}$-কে !!{derive}s
করে; আর $\Th{T}$ যেহেতু $\Th{Q}$-কে প্রসারিত করে, $\Th{T}$-ও তাই করে।
আমরা দেখালাম, $\Th{T}$ যদি $!G_\Th{T}$-কে !!{derive}s করে, তবে সেটি
$\lnot !G_\Th{T}$-কেও !!{derive}s করে; ফলে সেটি অসঙ্গতিপূর্ণ হতো।
\end{proof}

\begin{defn}
\ollabel{thm:oconsis-q}
তত্ত্ব $\Th{T}$-কে \emph{$\omega$-সঙ্গতিপূর্ণ} বলা হয় যদি নিচেরটি সত্য
হয়: $\lexists[x][!A(x)]$ যেকোনো বাক্য এবং $\Th{T}$ যদি $\lnot
!A(\num 0)$, $\lnot !A(\num 1)$, $\lnot !A(\num 2)$, \dots{}-কে
!!{derive}s করে, তবে $\Th{T}$ $\lexists[x][!A(x)]$-কে প্রমাণ করে না।
\end{defn}

খেয়াল করো, প্রতিটি $\omega$-সঙ্গতিপূর্ণ তত্ত্ব সঙ্গতিপূর্ণও। কারণ,
$\Th{T}$ অসঙ্গতিপূর্ণ হলে প্রতিটি~$!A$-এর জন্য $\Th{T} \Proves !A$।
বিশেষত, $\Th{T}$ অসঙ্গতিপূর্ণ হলে সেটি প্রতিটি~$n$-এর জন্য $\lnot
!A(\num n)$-কে !!{derive}s করে এবং $\lexists[x][!A(x)]$-কেও !!{derive}s
করে। অতএব $\Th{T}$ অসঙ্গতিপূর্ণ হলে সেটি $\omega$-অসঙ্গতিপূর্ণ। বিপরীত
প্রতিবচনে, $\Th{T}$ $\omega$-সঙ্গতিপূর্ণ হলে অবশ্যই সঙ্গতিপূর্ণ।

\begin{lem}\ollabel{lem:omega-cons-G-unref}
$\Th{T}$ যদি $\Th{Q}$-কে প্রসারিত করে এমন একটি $\omega$-সঙ্গতিপূর্ণ,
!!{axiomatizable} তত্ত্ব হয়, তবে $\Th{T} \Proves/ \lnot !G_\Th{T}$।
\end{lem}

\begin{proof}
দেখাব, $\Th{T}$ যদি $\lnot !G_\Th{T}$-কে !!{derive}s করে, তবে সেটি
$\omega$-অসঙ্গতিপূর্ণ। ধরি $\Th{T}$ সূত্র $\lnot !G_\Th{T}$-কে
!!{derive}s করে। $\Th{T}$ অসঙ্গতিপূর্ণ হলে সেটি $\omega$-অসঙ্গতিপূর্ণ,
তাই কাজ শেষ। নচেৎ $\Th{T}$ সঙ্গতিপূর্ণ; ফলে
\olref{lem:cons-G-unprov} অনুসারে সেটি $!G_\Th{T}$-কে !!{derive} করে না।
$\Th{T}$-তে $!G_\Th{T}$-এর কোনো !!{derivation} নেই বলে $\Th{Q}$
নিচেরগুলি !!{derive}s করে:
\[
\lnot \OPrf[\Th{T}](\num 0, \gn{!G_\Th{T}}), \lnot \OPrf[\Th{T}](\num 1,
\gn{!G_\Th{T}}), \lnot \OPrf[\Th{T}](\num 2, \gn{!G_\Th{T}}), \dots
\]
এবং $\Th{T}$-ও তাই করে। অন্যদিকে \olref{eqn:qpf} অনুসারে $\lnot
!G_\Th{T}$ হলো
$\lexists[x][\OPrf[\Th{T}](x,\gn{!G_\Th{T}})]$-এর সমতুল্য। অতএব
$\Th{T}$ $\omega$-অসঙ্গতিপূর্ণ।
\end{proof}

\begin{prob}
  প্রতিটি $\omega$-সঙ্গতিপূর্ণ তত্ত্ব সঙ্গতিপূর্ণ। দেখাও যে বিপরীতটি সত্য
  নয়, অর্থাৎ সঙ্গতিপূর্ণ কিন্তু $\omega$-অসঙ্গতিপূর্ণ তত্ত্ব আছে।
  $\Th{Q} \cup \{\lnot !G_\Th{Q}\}$ সঙ্গতিপূর্ণ কিন্তু
  $\omega$-অসঙ্গতিপূর্ণ দেখিয়ে এটি প্রমাণ করো।
\end{prob}

\begin{thm}
\ollabel{thm:first-incompleteness} ধরি $\Th{T}$ হলো $\Th{Q}$-কে
প্রসারিত করে এমন যেকোনো $\omega$-সঙ্গতিপূর্ণ, !!{axiomatizable} তত্ত্ব।
তবে $\Th{T}$ সম্পূর্ণ নয়।
\end{thm}

\begin{proof}
  $\Th{T}$ $\omega$-সঙ্গতিপূর্ণ হলে সেটি সঙ্গতিপূর্ণ; ফলে
  \olref{lem:cons-G-unprov} অনুসারে $\Th{T} \Proves/ !G_\Th{T}$।
  \olref{lem:omega-cons-G-unref} অনুসারে $\Th{T} \Proves/ \lnot
  !G_\Th{T}$। অর্থাৎ $\Th{T}$ অসম্পূর্ণ, কারণ সেটি $!G_\Th{T}$ ও
  $\lnot !G_\Th{T}$—কোনোটিই !!{derive}s করে না।
\end{proof}

\end{document}
